\documentclass[a4paper]{article}
\usepackage{amsmath}
\usepackage{amssymb}
\usepackage{geometry}
\usepackage{hyperref}

\geometry{margin=25mm}

\title{The Log Chart of the Three-Sphere\\
and the Intermediate-Axis Reversal}

\author{https://github.com/hypernumbernet}
\date{\today}

\begin{document}

\maketitle

\begin{abstract}
Particles sit on the three-sphere of unit quaternions, and the relative position of two of them is the rotor \(\rho(a,b)=a^{-1}b\). The force follows a great circle. The laboratory reading is the azimuthal equidistant chart, \(\log(\cos\chi+n\sin\chi)=\chi n\) for \(0<\chi<\pi\). The chart is faithful at the origin and singular at the antipode. A rotation about the third chart axis, and the reflection through the plane of the first two, commute with the chart. Left multiplication by \(\cos t+I\sin t\) leaves every relative rotor unchanged. On the point whose chart position is \((0,\chi,0)\), the same multiplication has chart velocity \((\chi\cot\chi,\,0,\,\chi)\) at the identity instant, the sum of a term along the first axis and the cross product of that axis with the chart position. At the equator the first term vanishes and the velocity is the cross product. The moment of one pushed great circle about another body vanishes on a common ray and when the source is the chart origin. At equal chart distance it equals \((1-\cos\alpha)\,\chi\,(\chi-\sin\chi\cos\chi)\sin\alpha\), positive for angles in \((0,\pi)\), and odd in the angle, so a mirror cancels it. Six unit masses at the intercepts \((\pm a,0,0)\), \((0,\pm b,0)\), \((0,0,\pm c)\) therefore exert no chart torque: in each coordinate plane the four right-angle moments, two of each sign, sum to zero. The products of inertia vanish. The principal moments are \(I_x=2(b^2+c^2)\), \(I_y=2(a^2+c^2)\), \(I_z=2(a^2+b^2)\), so \(0<a<b<c\) gives \(I_z<I_y<I_x\), and two moments agree if and only if the two corresponding lengths agree. One partner moved off the ray by an angle in \((0,\pi)\), at equal chart distance, leaves a positive moment, and the mirror of that partner cancels it. A cloud at rest thereby acquires angular momentum. The rotation is the chart reading of the spherical forces. When the torque vanishes and one rotation carries the shape, the body angular velocity obeys Euler's equations, and of the three principal spins only the intermediate axis is linearly unstable. For \(0<I_1<I_2<I_3\) and \(B>0\),
\[
\omega_1=A\operatorname{sech}(kt),\qquad
\omega_2=B\tanh(kt),\qquad
\omega_3=C\operatorname{sech}(kt),
\]
with \(k\), \(A\), and \(C\) fixed by the moments and by \(B\), satisfies those equations, keeps the energy \(\tfrac12 I_2 B^2\), and keeps \(L^2=2EI_2\). The intermediate component changes sign once. A residual of relative size \(\varepsilon<1/6\), growing at this rate, takes longer to reach half amplitude than the reversal takes to pass between the two half-amplitude states, and that wait exceeds every bound as the residual shrinks, while the crossing scale \(1/k\) stays fixed. The aggregate enters the free motion only through its inertia tensor. Over the \(2^n\) assignments of a sign to each of \(n\) moments, the uniform average of the signed sum is zero, and the uniform average of its square is the sum of the squares. The canonical weight at inverse temperature \(\beta>0\) is the product of three centered Gaussians of variances \(1/(\beta I_i)\), the normalized factor \(e^{-\beta E}\). The mean of the gap \(L^2-2EI_2\) is \((I_1+I_3-2I_2)/\beta\). That weight is positive on both signs of the gap, and on the transverse residuals of relative size in \((0,1/6)\). Each of those residuals waits longer than the reversal takes to cross.
\end{abstract}

\section{Introduction}

The double spacetime theory places each massive particle on a compact three-sphere~\cite{dst-main}. The relative position of two particles \(a\) and \(b\) is the rotor
\[
\rho(a,b)=a^{-1}b.
\]
It records where one particle sits as seen from the other. It is not an intra-particle mismatch. The force between them follows the great circle determined by that rotor. The laboratory chart is the azimuthal equidistant projection, the logarithm of the unit quaternion,
\begin{equation}
\label{eq:log}
\log(\cos\chi+n\sin\chi)=\chi n\in\mathbb{R}^3,
\qquad 0<\chi<\pi,
\end{equation}
with \(n\) a Euclidean unit vector. The chart is faithful at the origin and singular at the antipode, which the projection spreads along the boundary of the ball of radius \(\pi\). The volume distortion is the Jacobian \(\sin^2\chi/\chi^2\) of the main paper. The galactic rotation curve built there from the same chart is left in that paper. What is followed here is the torque that the chart assigns to a cloud, and the free rigid motion of that torque.

\section{The chart}
\label{sec:chart}

A rotation about the third chart axis fixes the pole, and so does the reflection through the plane of the first two axes. Both preserve the Euclidean norm of a direction, and both lift to the ambient space \(\mathbb{R}\times\mathbb{R}^3\) by acting on the imaginary part and fixing the real axis. On a lifted unit direction of radius \(\chi\in(0,\pi)\), the logarithm intertwines each of these maps with the same map of the chart. The chart of a rotated point is the rotation of the chart point, and the chart of a reflected point is the reflection of the chart point.

Left multiplication does not move relative positions. For every real \(t\) and every pair of unit quaternions,
\begin{equation}
\label{eq:left-relative}
\rho\bigl((\cos t+I\sin t)\,a,\;(\cos t+I\sin t)\,b\bigr)=\rho(a,b).
\end{equation}
The same multiplication moves the chart. Take the point \(\cos\chi+J\sin\chi\), whose chart position is \((0,\chi,0)\), and multiply on the left by \(\cos t+I\sin t\). At the instant \(t=0\) the chart velocity is
\begin{equation}
\label{eq:left-velocity}
\bigl(\chi\cot\chi,\;0,\;\chi\bigr)
=
\chi\cot\chi\,e_1
+
e_1\times(\chi e_2).
\end{equation}
The second summand is the cross product of the generator with the chart position. At the equator \(\chi=\pi/2\) the cotangent vanishes, and the velocity equals that cross product. For \(\chi\neq 0\) the third component \(\chi\) is not zero. A common left translation, invisible to every relative rotor, is a rotation in the chart.

\section{Two bodies}
\label{sec:two}

Push the initial velocity of the great circle into the chart. Its moment about the source is the chord moment of the main paper. The moment vanishes when \(\sin\alpha=0\), so a source on the same ray as the body contributes nothing, and it vanishes when the source is the chart origin, because the great circle through the origin is a straight ray of the chart. When the two chart distances are equal to \(\chi\in(0,\pi)\) and the oriented angle is \(\alpha\), the moment is
\begin{equation}
\label{eq:equal}
(1-\cos\alpha)\,\chi\,(\chi-\sin\chi\cos\chi)\sin\alpha.
\end{equation}
On \((0,\pi)\times(0,\pi)\) this quantity is positive. It is odd in \(\alpha\): the moment at \(-\alpha\) is the negative of the moment at \(\alpha\), and the two sum to zero. A mirror pair cancels.

At equal chart distance the aligned angle \(\alpha=0\) has moment zero, since \(\sin 0=0\). The same distance at an angle \(\alpha\in(0,\pi)\) has positive moment. Moving one partner off the common ray, or omitting the mirror that would have cancelled it, leaves a nonzero moment.

\section{A triaxial cloud}
\label{sec:cloud}

The chart torque of a cloud is the sum of these pairwise moments. Place six unit masses at the intercepts
\[
(\pm a,0,0),\qquad (0,\pm b,0),\qquad (0,0,\pm c),
\]
with the lengths inside the chart. Each pair lies in a coordinate plane. In one such plane the four right-angle moments, two of each sign, sum to zero, because each sign cancels the other. The three planes are the three components of the torque, and each component vanishes. The cloud of intercepts is torque-free.

The Euclidean second moments of the same six points are
\begin{equation}
\label{eq:inertia}
I_x=2(b^2+c^2),\qquad
I_y=2(a^2+c^2),\qquad
I_z=2(a^2+b^2).
\end{equation}
Every product of inertia vanishes, so these are the principal moments and the coordinate axes are the principal axes. If \(0<a<b<c\), then
\begin{equation}
\label{eq:order}
I_z<I_y<I_x.
\end{equation}
For nonnegative lengths, \(I_x=I_y\) if and only if \(a=b\), \(I_y=I_z\) if and only if \(b=c\), and \(I_z=I_x\) if and only if \(a=c\). Two principal moments agree only when two of the intercept lengths agree. That is a closed condition on the shape.

The torque-free cloud is the balanced one. Displace one partner, at equal chart distance, from the ray through the origin out to an angle \(\alpha\in(0,\pi)\). The aligned moment is zero and the displaced moment is positive, so the cloud, initially at rest, acquires angular momentum. Restoring the mirror partner at \(-\alpha\) cancels that moment again. The angular momentum is not supplied by an external couple. It is what remains when the spherical forces are read in the logarithm.

\section{Free rotation}
\label{sec:euler}

Suppose the chart torque remains zero and a single rotation carries the shape. The angular velocity in the body frame then obeys Euler's equations. With principal moments \(0<I_1<I_2<I_3\),
\begin{equation}
\label{eq:euler}
\begin{aligned}
\dot\omega_1&=\frac{I_2-I_3}{I_1}\,\omega_2\omega_3,\\
\dot\omega_2&=\frac{I_3-I_1}{I_2}\,\omega_3\omega_1,\\
\dot\omega_3&=\frac{I_1-I_2}{I_3}\,\omega_1\omega_2.
\end{aligned}
\end{equation}
A steady spin \(\Omega\neq 0\) along a principal axis linearises to a transverse rate whose square is fixed by the two gaps. About the intermediate axis that square is
\begin{equation}
\label{eq:rate-mid}
\frac{(I_3-I_2)(I_2-I_1)}{I_1 I_3}\,\Omega^2,
\end{equation}
which is positive. About the small axis the square is
\[
-\frac{(I_3-I_1)(I_2-I_1)}{I_2 I_3}\,\Omega^2,
\]
and about the large axis
\[
\frac{(I_2-I_3)(I_3-I_1)}{I_1 I_2}\,\Omega^2.
\]
Both are negative. A positive real number has positive square, so neither the small axis nor the large axis admits a positive real eigenvalue. The positive square root of~\eqref{eq:rate-mid} is a positive eigenvalue of the intermediate axis. If \(I_1=I_2\) or \(I_2=I_3\), the middle rate is zero. For the intercepts~\eqref{eq:inertia}, that degeneracy is the agreement of two lengths.

\section{The reversal}
\label{sec:reversal}

The unstable axis carries an explicit orbit. Let \(B>0\) and set
\begin{equation}
\label{eq:amplitudes}
\begin{aligned}
k&=B\sqrt{\frac{(I_2-I_1)(I_3-I_2)}{I_1 I_3}},\\[4pt]
A&=B\sqrt{\frac{I_2(I_3-I_2)}{I_1(I_3-I_1)}},\\[4pt]
C&=B\sqrt{\frac{I_2(I_2-I_1)}{I_3(I_3-I_1)}}.
\end{aligned}
\end{equation}
The profile
\begin{equation}
\label{eq:separatrix}
\omega_1=A\operatorname{sech}(kt),\qquad
\omega_2=B\tanh(kt),\qquad
\omega_3=C\operatorname{sech}(kt)
\end{equation}
satisfies~\eqref{eq:euler} at every \(t\). Its kinetic energy is the constant
\begin{equation}
\label{eq:energy}
E=\frac12 I_2 B^2,
\end{equation}
and the squared angular momentum is the constant \(I_2^2 B^2\), hence
\begin{equation}
\label{eq:critical}
L^2=2EI_2.
\end{equation}
The square of \(k\) is the intermediate rate~\eqref{eq:rate-mid} evaluated at \(\Omega=B\). The intermediate component vanishes if and only if \(t=0\). It is negative for \(t<0\) and positive for \(t>0\), so it changes sign once, and the orbit runs from \(-B\) to \(B\).

The half-amplitude states sit at a definite time. Since \(\tanh((\log 3)/2)=1/2\) and \(\tanh\) is odd,
\[
\omega_2=\pm\frac B2
\quad\text{at}\quad
t=\pm\frac{\log 3}{2k}.
\]
The transit from one to the other occupies
\begin{equation}
\label{eq:crossing}
\frac{\log 3}{k}.
\end{equation}
In the linear growth about the intermediate axis, a residual of relative amplitude \(\varepsilon\in(0,1/2)\) reaches half amplitude after
\begin{equation}
\label{eq:waiting}
\frac1k\log\frac{1}{2\varepsilon}.
\end{equation}
Whenever \(0<\varepsilon<1/6\), the logarithm \(\log(1/(2\varepsilon))\) exceeds \(\log 3\), and the wait~\eqref{eq:waiting} exceeds the transit~\eqref{eq:crossing}. The same wait exceeds every bound: for each real \(N\) there is a positive \(\varepsilon<1/6\) whose wait is larger than \(N\). The crossing scale \(1/k\) does not depend on \(\varepsilon\). An orbit close to~\eqref{eq:critical} may therefore linger, and then reverse on the scale already fixed by the moments and by \(B\).

The cloud enters this motion only as the triple \((I_1,I_2,I_3)\). For the intercepts, that triple is~\eqref{eq:inertia}, and an equal pair of axes is the agreement of two of \(a\), \(b\), and \(c\). A defect that leaves a positive moment has a direction. The component of that residual along the intermediate axis is the unstable one, and~\eqref{eq:separatrix} is the orbit in which that component reverses. This is the intermediate-axis reversal of a free triaxial body~\cite{ashbaugh1991}. The comparison of~\eqref{eq:waiting} with~\eqref{eq:crossing} is the linear growth against the separatrix transit. It is not the period of a neighbouring orbit.

\section{A mirror ensemble}
\label{sec:mirror}

The equal-distance moment~\eqref{eq:equal} is odd in the angle, so a partner and its mirror cancel. Assign each of \(n\) partners a sign \(\sigma_j=\pm 1\), and let \(m_j\) be the moment that partner contributes. The uniform average of \(\sum_j\sigma_j m_j\), taken over the \(2^n\) assignments of the signs, is zero. The uniform average of the square of that sum is \(\sum_j m_j^2\). If some \(m_j\) is nonzero, the average of the square is positive. A mirror-symmetric cloud has no preferred sense of rotation, and a positive mean-square chart torque. The angular momentum read by the logarithm is this mean square.

\section{The canonical weight}
\label{sec:canonical}

On the torque-free motion the energy is \(E=\tfrac12\sum_i I_i\omega_i^2\) and \(L^2=\sum_i I_i^2\omega_i^2\). The gap against~\eqref{eq:critical} is
\begin{equation}
\label{eq:gap}
L^2-2EI_2=I_1(I_1-I_2)\omega_1^2+I_3(I_3-I_2)\omega_3^2.
\end{equation}
The intermediate component does not appear. The gap vanishes at every spin \((0,\omega_2,0)\). For \(0<I_1<I_2<I_3\) it is negative at \((1,0,0)\) and positive at \((0,0,1)\).

The canonical weight at inverse temperature \(\beta>0\) is the product of three centered Gaussians, of variances \(1/(\beta I_i)\). That product is the normalized factor \(e^{-\beta E}\). Each principal term \(\tfrac12 I_i\omega_i^2\) has mean \(1/(2\beta)\), and the mean energy is \(3/(2\beta)\). The mean of the gap is
\begin{equation}
\label{eq:mean-gap}
\frac{I_1+I_3-2I_2}{\beta}.
\end{equation}
It vanishes if and only if \(I_2\) is the arithmetic mean of \(I_1\) and \(I_3\). For the intercepts~\eqref{eq:inertia} that condition is \(2b^2=a^2+c^2\). Among \(0<a<b<c\) both signs occur: the lengths \(1\), \(2\), \(3\) make the mean negative, and the lengths \(1\), \(2\), \(5/2\) make it positive.

The same weight is positive on every ball. About an intermediate spin \((0,B,0)\) the transverse residual
\[
\varepsilon=\sqrt{\frac{I_1\omega_1^2+I_3\omega_3^2}{I_2 B^2}}
\]
vanishes on the axis itself. The state \(\omega_1=B\sqrt{I_2/I_1}/12\), \(\omega_3=0\) has residual \(1/12\). The set of states whose residual lies in \((0,1/6)\) is open and contains that state, so it contains a ball about the state, and the weight of the set is positive. Every residual in the set is below one sixth, so the wait~\eqref{eq:waiting} exceeds the transit~\eqref{eq:crossing}. The canonical weight therefore charges the reversals that linger and then change the sign of the intermediate component.

\section{Conclusion}

The logarithm~\eqref{eq:log} is faithful at the origin, singular at the antipode, and equivariant under the pole-fixing rotation and reflection of~\S\ref{sec:chart}. A common left translation preserves every relative rotor~\eqref{eq:left-relative} and contributes the cross product in the chart velocity~\eqref{eq:left-velocity}. Pairwise, the chord moment vanishes on a ray and at the origin, is given at equal distance by~\eqref{eq:equal}, and cancels against its mirror. The six intercepts are torque-free, with principal moments~\eqref{eq:inertia} ordered by~\eqref{eq:order} when \(a<b<c\). One partner off the ray leaves a positive moment, and a cloud at rest acquires angular momentum from the spherical forces as the chart reads them.

On the torque-free rigid motion the body angular velocity satisfies~\eqref{eq:euler}. Only the intermediate axis has a positive real linear eigenvalue. The profile~\eqref{eq:separatrix} realises that instability, holds the energy~\eqref{eq:energy} and the critical value~\eqref{eq:critical}, and reverses its intermediate component once. A residual below one sixth of the amplitude waits longer than the reversal takes to cross, and the wait lengthens without bound as the residual shrinks.

Over the sign assignments of~\S\ref{sec:mirror} the uniform average of the signed sum of chart moments is zero, and the average of its square is the sum of the squares. The canonical weight of~\S\ref{sec:canonical} is the product of the three Gaussians of variances \(1/(\beta I_i)\), the normalized factor \(e^{-\beta E}\). Its mean gap is~\eqref{eq:mean-gap}. The weight is positive on both signs of~\eqref{eq:gap}, and on the residuals in \((0,1/6)\), each of which waits longer than the crossing.

The galactic curve of the main paper is another use of the chord moment. The rotation here is the torque of the triaxial cloud.

No floating-point integration of a many-body trajectory is claimed. No frequency of reversals is drawn from the weight. No value of \(a_0\) is derived, and the cloud is not identified with a galaxy.

\begin{thebibliography}{9}
\raggedright

\bibitem{dst-main}
Anonymous,
``Gravity as Torsion between Dual Spacetime: A Biquaternionic Reformulation of General Relativity'' (2025).
\url{https://github.com/hypernumbernet/dual-spacetime-doc/blob/main/double-spacetime-theory.pdf}.

\bibitem{ashbaugh1991}
M.~S. Ashbaugh, C.~C. Chicone, and R.~H. Cushman,
``The twisting tennis racket,''
\textit{J.\ Dynam.\ Differential Equations} \textbf{3}, 67--85 (1991).

\end{thebibliography}

\end{document}
